\scanpage{22}

\textcircled{15}

On fixe entier $n$ premier aux
car. résiduelles.

\uline{Classe fondamentale locale}

Supposons $Y\longrightarrow X$ ``intersection complète de
dim relative $1$'' et supposons d'abord
$Y=V(\varphi)$, $\varphi$ section régulière de $\mathcal O_X$. Alors
si $u:X-Y=U\longrightarrow X$, donc
$\deleted{i_*(\mathcal O_U)\simeq\mathcal O_X[1/\varphi]}$,
$\deleted{i_*(\mathcal O_U^*)\longrightarrow}$
$\varphi\in\Gamma(i_*(G_{mU}))$
$\deleted{\unsure{et}\ \unsure{\sim}\ \in\Gamma(\unsure{X,i_*(G_{mU}^*)})}$.
\deleted{D'où} alors $\varphi$ est déterminé \inserted{par $Y$}
(mod multiplication
par un $\in\Gamma(X,\unsure{i_*}(G_{mX}))$), donc \deleted{l'image} $\varphi$ définit
une section \deleted{$\varphi\in$} \uline{canonique} de
$\deleted{i_*(\mathcal O_U^*)/}$
$i_*(G_{mU})/G_{mX}$. On a la suite exacte
canonique
\diagram{\begin{tikzcd}[column sep=small]
0\arrow[r]&\underline H_Y^0(G_{mX})\arrow[r]\arrow[d,equal]&
G_{mX}\arrow[r,dash]&i_*(G_{mU})\arrow[r,dash,"d"]&
\underline H_Y^1(G_{mX})\arrow[d]\\
&0&&&0
\end{tikzcd}}
qui donne ici
\[
\deleted{\unsure{H^1}}\ \underline H_Y^1(G_{mX})
\xrightarrow{\sim}i_*(G_{mU})/\operatorname{Im}G_{mX}
\]
et montre que $d\varphi$ est une section
\uline{canonique} de $\underline H_Y^1(G_{mX})$, que nous allons
noter $\deleted{\unsure{\varepsilon_{Y/X}}}\ \unsure{\beta_{Y/X}}$.
D'autre part,
on a \uline{la} suite exacte de Kummer
\diagram{\begin{tikzcd}
0\arrow[r,dash]&\mu_{nX}\arrow[r]&G_{mX}\arrow[r,dash,"n"]&G_{mX}\arrow[r]&0
\end{tikzcd}}

\mspage{21}
